% GENERATED FILE. Edit canonical lessons, then run npm run generate:books.
% canonical-source-hash: fnv1a64:9fd1ecb03c4cacf7
% canonical-lessons: PT-C365-justica, PT-C365-injusto, PT-C365-injustica, PT-C365-preconceito, PT-C365-estereotipo

\chapter{Justo e injusto}
\label{ch:pt-justo-e-injusto}
\hlchaptermodality{365}

\begin{chapteropening}
\textbf{By the end of this chapter:} \emph{I can talk about justice and injustice, prejudice and stereotypes.}

Complete the last lesson of chapter 365: I can talk about justice and injustice, prejudice and stereotypes.
\end{chapteropening}

\section[a justiça]{a justiça — the justice; the courts}

\label{lesson:PT-C365-justica}



\begin{quote}
Before the new one: say the Portuguese for to show, then the Portuguese for to demand.
\end{quote}

\subsection*{You'll want to know: a justiça}
\textbf{a justiça} — \textquotedblleft{}the justice; the courts\textquotedblright{}.

\begin{quote}
As pessoas querem justiça.
\end{quote}

\begin{quote}
\emph{People want justice.}
\end{quote}

\emph{Fazer justiça} — to do justice. \emph{A justiça} also means the courts as a system.

\begin{cousinweb}
\emph{justiça} comes from Latin \emph{iūstitia} \textquotedblleft{}justice, equity\textquotedblright{} (Lewis \& Short), through Old Portuguese \emph{justiça} (EtymDB); it is built on \emph{iūstus} \textquotedblleft{}fair, righteous\textquotedblright{}.
\end{cousinweb}

\begin{grammarlens}[title={What you've built}]
One more word for telling, explaining and describing at B1.
\end{grammarlens}

\subsection*{Your turn}
\begin{itemize}
\raggedright
  \item \emph{Say it:} \emph{a justiça}
  \item \emph{Say it:} \emph{a justiça}, once more
  \item \emph{Say it:} a sentence of your own with \emph{a justiça}
  \item \emph{Recall:} say the Portuguese for to show, then the Portuguese for to demand, then say \emph{a justiça} again
\end{itemize}

\begin{tcolorbox}[breakable,colback=teal!4,colframe=teal!35!black,title={Before you move on}]
What does a justiça mean? (\textquotedblleft{}The justice; the courts\textquotedblright{}.) Say it once more.
\end{tcolorbox}
\section[injusto]{injusto — unfair, unjust}

\label{lesson:PT-C365-injusto}



\begin{quote}
Before the new one: say the Portuguese for to demand, then the Portuguese for the justice.
\end{quote}

\subsection*{You'll want to know: injusto}
\textbf{injusto} — \textquotedblleft{}unfair, unjust\textquotedblright{}.

\begin{quote}
Acho injusto pagar mais pelo mesmo trabalho.
\end{quote}

\begin{quote}
\emph{I think it is unfair to pay more for the same work.}
\end{quote}

Feminine \emph{injusta}. The opposite of \emph{justo}.

\begin{cousinweb}
\emph{injusto} is a learned borrowing from Latin \emph{iniūstus} \textquotedblleft{}unjust\textquotedblright{} (Lewis \& Short), built from \emph{in-} \textquotedblleft{}not\textquotedblright{} and \emph{iūstus} \textquotedblleft{}fair, righteous\textquotedblright{}.
\end{cousinweb}

\begin{grammarlens}[title={What you've built}]
One more word for telling, explaining and describing at B1.
\end{grammarlens}

\subsection*{Your turn}
\begin{itemize}
\raggedright
  \item \emph{Say it:} \emph{injusto}
  \item \emph{Say it:} \emph{injusto}, once more
  \item \emph{Say it:} a sentence of your own with \emph{injusto}
  \item \emph{Recall:} say the Portuguese for to demand, then the Portuguese for the justice, then say \emph{injusto} again
\end{itemize}

\begin{tcolorbox}[breakable,colback=teal!4,colframe=teal!35!black,title={Before you move on}]
What does injusto mean? (\textquotedblleft{}Unfair, unjust\textquotedblright{}.) Say it once more.
\end{tcolorbox}
\section[a injustiça]{a injustiça — the injustice}

\label{lesson:PT-C365-injustica}



\begin{quote}
Before the new one: say the Portuguese for the justice, then the Portuguese for unfair.
\end{quote}

\subsection*{You'll want to know: a injustiça}
\textbf{a injustiça} — \textquotedblleft{}the injustice\textquotedblright{}.

\begin{quote}
Ela luta contra a injustiça.
\end{quote}

\begin{quote}
\emph{She fights against injustice.}
\end{quote}

\emph{Que injustiça!} — how unfair!

\begin{cousinweb}
\emph{injustiça} is a learned borrowing from Latin \emph{iniūstitia} \textquotedblleft{}injustice\textquotedblright{} (EtymDB), from \emph{iniūstus} \textquotedblleft{}unjust\textquotedblright{}, built on \emph{iūstus} \textquotedblleft{}fair, righteous\textquotedblright{} (Lewis \& Short).
\end{cousinweb}

\begin{grammarlens}[title={What you've built}]
One more word for telling, explaining and describing at B1.
\end{grammarlens}

\subsection*{Your turn}
\begin{itemize}
\raggedright
  \item \emph{Say it:} \emph{a injustiça}
  \item \emph{Say it:} \emph{a injustiça}, once more
  \item \emph{Say it:} a sentence of your own with \emph{a injustiça}
  \item \emph{Recall:} say the Portuguese for the justice, then the Portuguese for unfair, then say \emph{a injustiça} again
\end{itemize}

\begin{tcolorbox}[breakable,colback=teal!4,colframe=teal!35!black,title={Before you move on}]
What does a injustiça mean? (\textquotedblleft{}The injustice\textquotedblright{}.) Say it once more.
\end{tcolorbox}
\section[o preconceito]{o preconceito — the prejudice}

\label{lesson:PT-C365-preconceito}



\begin{quote}
Before the new one: say the Portuguese for unfair, then the Portuguese for the injustice.
\end{quote}

\subsection*{You'll want to know: o preconceito}
\textbf{o preconceito} — \textquotedblleft{}the prejudice\textquotedblright{}.

\begin{quote}
Ainda há muito preconceito contra os idosos.
\end{quote}

\begin{quote}
\emph{There is still a lot of prejudice against old people.}
\end{quote}

\emph{Ter preconceitos} — to be prejudiced.

\begin{cousinweb}
\emph{preconceito} is built from \emph{pre-} \textquotedblleft{}before\textquotedblright{} and \emph{conceito} \textquotedblleft{}idea\textquotedblright{}, from Latin \emph{conceptus} \textquotedblleft{}a collecting, a conceiving\textquotedblright{}, on \emph{concipere} \textquotedblleft{}to take in\textquotedblright{}, built on \emph{capere} \textquotedblleft{}to take\textquotedblright{} (Lewis \& Short): an idea formed before the facts.
\end{cousinweb}

\begin{grammarlens}[title={What you've built}]
One more word for telling, explaining and describing at B1.
\end{grammarlens}

\subsection*{Your turn}
\begin{itemize}
\raggedright
  \item \emph{Say it:} \emph{o preconceito}
  \item \emph{Say it:} \emph{o preconceito}, once more
  \item \emph{Say it:} a sentence of your own with \emph{o preconceito}
  \item \emph{Recall:} say the Portuguese for unfair, then the Portuguese for the injustice, then say \emph{o preconceito} again
\end{itemize}

\begin{tcolorbox}[breakable,colback=teal!4,colframe=teal!35!black,title={Before you move on}]
What does o preconceito mean? (\textquotedblleft{}The prejudice\textquotedblright{}.) Say it once more.
\end{tcolorbox}
\section[o estereótipo]{o estereótipo — the stereotype}

\label{lesson:PT-C365-estereotipo}



\begin{quote}
Before the new one: say the Portuguese for the injustice, then the Portuguese for the prejudice.
\end{quote}

\subsection*{You'll want to know: o estereótipo}
\textbf{o estereótipo} — \textquotedblleft{}the stereotype\textquotedblright{}.

\begin{quote}
Muitos filmes repetem o mesmo estereótipo.
\end{quote}

\begin{quote}
\emph{Many films repeat the same stereotype.}
\end{quote}

First a printing plate cast in one solid piece.

\begin{cousinweb}
\emph{estereótipo} is built from Greek \emph{stereós} \textquotedblleft{}solid\textquotedblright{} and \emph{tipo} \textquotedblleft{}type\textquotedblright{} (EtymDB), probably through French \emph{stéréotype}; it first named a printing plate cast in one solid piece.
\end{cousinweb}

\begin{grammarlens}[title={What you've built}]
One more word for telling, explaining and describing at B1.
\end{grammarlens}

\subsection*{Your turn}
\begin{itemize}
\raggedright
  \item \emph{Say it:} \emph{o estereótipo}
  \item \emph{Say it:} \emph{o estereótipo}, once more
  \item \emph{Say it:} a sentence of your own with \emph{o estereótipo}
  \item \emph{Recall:} say the Portuguese for the injustice, then the Portuguese for the prejudice, then say \emph{o estereótipo} again
\end{itemize}

\begin{tcolorbox}[breakable,colback=teal!4,colframe=teal!35!black,title={Before you move on}]
What does o estereótipo mean? (\textquotedblleft{}The stereotype\textquotedblright{}.) Say it once more.
\end{tcolorbox}
